%%%%%
%%%%%  Naudokite LUALATEX, ne LATEX.
%%%%%
%%%%
\documentclass[]{VUMIFTemplateClass}

\usepackage{indentfirst}
\usepackage{amsmath, amsthm, amssymb, amsfonts}
\usepackage{mathtools}
\usepackage{physics}
\usepackage{graphicx}
\usepackage{verbatim}
\usepackage[hidelinks]{hyperref}
\usepackage{color,algorithm,algorithmic}
\usepackage[nottoc]{tocbibind}
\usepackage{tocloft}

\usepackage{titlesec}
\newcommand{\sectionbreak}{\clearpage}

\makeatletter
\renewcommand{\fnum@algorithm}{\thealgorithm}
\makeatother
\renewcommand\thealgorithm{\arabic{algorithm} algoritmas}

\usepackage{biblatex}
\bibliography{bibliografija}
%% norint pakeisti bibliografijos šaltinių numeravimą (skaitiniu arba raidiniu), pakeitimus atlikti VUMIFTemplateClass.cls 150 eilutėje

% Author's MACROS
\newcommand{\EE}{\mathbb{E}\,} % Mean
\newcommand{\ee}{{\mathrm e}}  % nice exponent
\newcommand{\RR}{\mathbb{R}}


\studijuprograma{Programų sistemų} %Studijų programą įrašyti kilmininko linksniu (pavyzdžiui – Programų sistemų, Finansų ir draudimų matematikos ir t. t.)
\darbotipas{Bakalauro baigiamojo darbo planas} % Bakalauro baigiamasis darbas arba magistro baigiamasis darbas
\darbopavadinimas{Varžymosi principais grįstų atakų aptikimas naudojant paaiškinamo dirbtinio intelekto metodą kenkėjiškų programų kontekste}
\darbopavadinimasantras{Defense Against Adversarial Malware Obfuscation Attacks Using Explainable Artificial Intelligence}
\autorius{Liudas Kasperavičius}

%Autorių gali būti ir daugiau, tuo atveju, kiekvienas autorius rašomas iš naujos eilutės, ir pridedamas titulinis.tex arba dvigubasTitulinis.tex dokumentuose
%\antrasautorius{Vardas Pavardė} %Jei toks yra, kitu atveju ištrinti

\vadovas{prof. dr. Olga Kurasova}
% \recenzentas{pedagoginis/mokslinis vardas Vardas Pavardė} %Jei toks yra žinomas, kitu atveju ištrinti
% \moksliniskonsultantas{pedagoginis/mokslinis vardas Vardas Pavardė} %Jei toks yra žinomas, kitu atveju ištrinti

\begin{document}
\selectlanguage{lithuanian}

\onehalfspacing
\begin{titlepage}
\vskip 20pt
\begin{center}
\includegraphics[scale=0.55]{images/MIF.png}
\end{center}

\makeatletter

\vskip 20pt
\centerline{\bf \large \textbf{VILNIAUS UNIVERSITETAS}}
\vskip 10pt
\centerline{\large \textbf{MATEMATIKOS IR INFORMATIKOS FAKULTETAS}}
\vskip 10pt
\centerline{\large \textbf{\MakeUppercase{\@studijuprograma \space studijų programa}}}

\vskip 80pt
\centerline{\Large \@darbotipas}
\vskip 20pt
\begin{center}
    {\bf \LARGE \@darbopavadinimas}
\end{center}
\begin{center}
    {\bf \Large \@darbopavadinimasantras}
\end{center}
\vskip 80pt

\centering{\Large \@autorius}
\@ifundefined{@antrasautorius}{}
{
\vskip 10pt
\centering{\Large \@antrasautorius}
}
\vskip 20pt

\centering{
    \begin{tabular}{rcp{.7\textwidth}}
        {\Large Darbo vadovas} & {\Large :} & {\Large \@vadovas}\\[10pt]
        \@ifundefined{@moksliniskonsultantas}{}
            {
                {\Large Mokslinis konsultantas} & {\Large :} & {\Large \@moksliniskonsultantas}\\[10pt]
            }
        \@ifundefined{@recenzentas}{}
            {
                {\Large Recenzentas} & {\Large :} & {\Large \@recenzentas}\\[10pt]
            }
    \end{tabular}}


\vskip 110pt

\centerline{\large \textbf{Vilnius}}
\centerline{\large \textbf{\the\year{}}}

\makeatother

\newpage
\end{titlepage}
%\newgeometry{top=2cm,bottom=2cm,right=2cm,left=3cm}
\setcounter{page}{2}


\section{Darbo planas}

\subsection{Tyrimo objektas ir aktualumas}

Pagrindinis tyrimo objektas -- varžymosi principais grįstos atakos prieš
kenkėjiškų programų detektorius. Varžymosi principais grįstos atakos -- tai
subtilus duomenų modifikacijos metodas. Šių atakų metu duomenys pakeičiami
nežymiai, tačiau tokiu būdu, jog geba suklaidinti jiems pritaikytus dirbtiniu
intelektu įgyvendintus klasifikatorius. Vienas iš tokių atakų panaudojimo
atvejų -- kenkėjiškų programų obfuskacija -- procesas, kurio metu kenkėjiškos
programos pakeičiamos taip, jog jų veikimas išlieka toks pat, tačiau kenkėjiškų
programų detektoriai šių modifikuotų programų neaptinka (klasifikuoja kaip
nekenksmingas). Kadangi dauguma kenkėjiškų programų detektorių yra įgyvendinti
kaip dirbtinio intelekto algoritmai ar modeliai, svarbu ieškoti būdu aptikti
tokias atakas. Vienas iš galimų ir mokslinėje literatūroje minimų būdų
varžymosi principais grįstų atakų aptikimui -- dirbtinio intelekto modelio
priimtų sprendimų paaiškinimas ir jo analizė. Kadangi šiuolaikiniai dirbtinio
intelekto modeliai yra sudėtingi, daugiadimensiai ir sunkiai ar visiškai
neinterpretuojami, modelio sprendimų priėmimo paaiškinimas suprantamas kaip
paprastesnio (klasikinio statistinio) modelio ar metodo panaudojimas
suprantamesnei atskiro atvejo klasifikacijos ir klasifikavimo proceso
interpretacijai.

\subsection{Darbo tikslas}

Sukurti metodą sėkmingų varžymosi principais grįstų atakų aptikimui ir
paaiškinimui.

\subsection{Keliami uždaviniai ir laukiami rezultatai}

\subsubsection{Uždaviniai}
\begin{enumerate}
    \item Apibrėžti varžymosi principais grįstų atakų aptikimo metodą naudojant LIME
          (\textit{angl. Local Interpretable Model-Agnostic Explanations}).
    \item Pritaikyti varžymosi principais grįstų atakų aptikimą dvejetainius požymių
          vektorius naudojantiems modeliams taikant dimensijų mažinimo metodus.
    \item Sukurti klasifikavimo proceso praplėtimą į jį įtraukiant varžymosi principais
          grįstų atakų aptikimą.
    \item Ištirti praplėsto klasifikavimo proceso tikslumą.
\end{enumerate}

\subsubsection{Laukiami rezultatai}
\begin{enumerate}
    \item Sukurtas varžymosi principais pagrįstų atakų aptikimo metodas, pritaikomas bet
          kokiam kenkėjiškų programų detektoriui.
    \item Laikantis klasifikavimo proceso praplėtimo aptinkama daugiau obfuskuotų
          kenkėjiškų programų, nei jo nesilaikant.
    \item Klasifikavimo proceso praplėtimas paaiškinamo dirbtinio intelekto metodu padeda
          nustatyti pažeidžiamus požymius (modelio silpnąsias vietas)
\end{enumerate}

\subsection{Tyrimo metodai}

\begin{enumerate}
    \item \textbf{Literatūros analizė}.
    \item \textbf{Modeliavimo metodas}. Klasifikavimo proceso praplėtimo tyrime bus naudojami šie modeliai:
          \begin{itemize}
              \item kenkėjiškų programų klasifikatoriaus modelis, įgyvendintas kaip mašininio
                    mokymosi algoritmas,
              \item varžymosi principais grįstų atakų generavimo modelis, įgyvendintas kaip
                    \textit{MalGAN} karkasą atitinkantis mašinino mokymosi modelis.
          \end{itemize}
    \item \textbf{Lyginamosios analizės metodas}. Lyginami klasifikavimo (kenkėjiškų programų aptikimo) rezultatai taikant klasifikavimo proceso praplėtimą ir jo netaikant.
\end{enumerate}

\subsection{Numatomas darbo atlikimo procesas}

\begin{enumerate}
    \item Literatūros analizė.
    \item Kenkėjiškų ir nekenkėjiškų programų požymių duomenų rinkinio pasirinkimas.
    \item Varžymosi principais grįstų atakų karkaso ir kenkėjiškų programų detektoriaus
          modelio mokymas.
    \item \textit{LIME} komponento paruošimas varžymosi principais grįstų atakų aptikimui ir paaiškinimui (\textit{normalių požymių} aibės kūrimas iš mokymosi duomenų rinkinio).
    \item Dvejetainių požymių vektorių supaprastinimas mažinant dimensijas ir šio
          komponento integracija su \textit{LIME} komponentu.
    \item Klasifikavimo proceso praplėtimo apibrėžimas ir iliustracijų rengimas.
    \item Klasifikavimo proceso praplėtimo tikslumo tyrimas taikant lyginamąją analizę.
\end{enumerate}

\subsection{Literatūros šaltinių apibūdinimas}

Mokslinėje literatūroje aptariami varžymosi principais grįstų atakų kūrimo
karkasai. Vienas iš populiariausių -- \textit{MalGAN}
karkasas~\cite{huGeneratingAdversarialMalware2017} dažnai naudojamas kaip
bazinis tolimesniems tyrimams. Apie paaiškinamo dirbtinio intelekto
(specifiškai \textit{LIME}~\cite{ribeiroWhyShouldTrust2016}) naudojimą įsilaužimo aptikimo sistemose
(\textit{angl. Intrusion Detection Systems, IDS}) siekiant aptikti varžymosi
principais grįstas atakas kalba Tcydenova et
al.~\cite{tcydenovaDetectionAdversarialAttacks2021}. \textit{IDS} yra realiu
laiku veikiančios ir dažniausiai dinaminius požymius analizuojančios sistemos,
tačiau šis tyrimas apsiriboja statinių požymių, kurių apibendrinimą atliko
Demetrio et al.~\cite{demetrioAdversarialEXEmplesSurvey2021}, analize.
Potencialūs statinių požymių rinkiniai, kurie gali būti tinkami šiam tyrimui
yra \textit{EMBER}~\cite{andersonEMBEROpenDataset2018a} ir
\textit{SLEIPNIR}~\cite{al-dujailiAdversarialDeepLearning2018}.

\end{document}
